\section{Related literature and the contribution boundary}
\subsection{Structural change: augmentation, displacement and persistent scarcity}
Task-based accounts distinguish labour-displacing automation from productivity and new-task effects \citep{acemoglu2018}. \citet{acemoglu2024} provides a restrained macroeconomic assessment, whereas \citet{trammell2026}, \citet{restrepo2025} and \citet{jones2026} consider more transformative possibilities and their limiting inputs. These are conditional theories, not alternative estimates of the same already-observed technological shock. Our declining labour-share schedule is a scenario assumption; it is not fitted to those papers or used to assert the disappearance of work.

This question also belongs to structural-change theory rather than only to monetary policy. \citet{wirkierman2023} connects demand, productivity, wages and public intervention in an empirically implemented structural accounting framework. \citet{gualerzi2026} emphasises demand and market creation in expansive structural change. The multi-sector model of \citet{dosi2022} shows that labour-market and distributive institutions can alter technological-unemployment outcomes. Our smaller laboratory does not reproduce their endogenous sectors or task dynamics; it isolates a participant-equity contract and its financing and coverage constraints.

The distinction matters empirically. \citet{brynjolfsson2025} study 5,172 customer-support agents and find a 15\% average productivity increase, with larger gains among less experienced workers. Their design does not identify economy-wide employment or wage responses. \citet{minniti2025} find heterogeneous associations between AI innovation and the labour share across European regions; patents and regional innovation are not measures of full production automation. \citet{katz2026}, in this journal, report productivity enhancement but modest observed aggregate effects during early generative-AI diffusion. This evidence motivates distributional monitoring and heterogeneous scenarios, not a forecast of an imminent post-labour economy. \citet{lamperti2026}, also in this journal, find that automation adoption is associated on average with a higher probability of inactivity in European labour-market data, with heterogeneous effects across demographic groups. At the same time, \citet{fannguyen2026} value the labour-time currently saved by observed AI usage at a labour-cost equivalent of USD~2.7 trillion, or 3.4\% of GDP; they explicitly treat that number as a valuation of time saved rather than realised GDP. Together these results strengthen the case for monitoring distributional channels while preserving uncertainty about the magnitude and timing of the macro transition.

Supply growth also depends on complementarity. Energy, compute infrastructure, chips, land and organisational capabilities can remain scarce. The simulation therefore caps potential output by a resource envelope. It does not turn the automation parameter into unlimited output, and its three activities do not constitute an input-output model of European supply chains. A useful structural-change contribution must explain both who receives productive gains and when productive constraints prevent nominal claims from becoming additional consumption.

\subsection{Taxes, insurance, industrial policy and the strongest alternatives}
The fiscal response to AI is not exhausted by a robot tax. \citet{brollo2024} discuss social protection, capital taxation and the broader fiscal response; \citet{korinek2026} analyse how transformative AI could change tax bases and the role of public capital. \citet{korinekstiglitz2026} distinguish redistribution from steering technical change. Existing capital and labour tax distortions can themselves favour excessive automation \citep{acemoglutax2020}. These contributions rule out a benchmark in which tax policy is held artificially incapable of adapting.

Optimal-tax work is particularly important as a counterargument. \citet{guerreiro2022} identify conditions for transitional robot taxation. \citet{thuemmel2023} finds that the optimal intervention can change from a subsidy to a tax and that, in the calibrated setting, most welfare gains arise through income-tax reform rather than a separate robot tax. These results do not settle the present institutional question, but they preclude claiming that an equity requirement is automatically less distortionary merely because it is not called a tax. Anticipated dilution reduces the value of an investor's future claim.

Subsidies can support learning, coordination, infrastructure and capacity rather than simply enrich current owners. \citet{juhasz2024} reviews evidence showing contextual and sometimes persistent industrial-policy effects. A subsidy can also be conditional on public or participant equity, and tax revenue can finance a social wealth fund. These are serious stock-building comparators. The present paper does not numerically optimise such an industrial-policy package; it identifies that omission rather than treating all subsidies as disposable consumption.

Cash transfers perform insurance and demand-support functions even without creating legal capital claims. Alaska's experience is informative about a permanent dividend in a market economy, not a national experiment in monetary financing of post-work incomes \citep{jonesmarinescu2020,bibler2023}. We use both targeted and universal transfers and keep their coverage differences explicit. A participation dividend cannot substitute for adequate support to a person unable to participate.

\subsection{Distributed capital, user contribution and corporate rights}
The lexical proposal in this paper has substantial conceptual ancestry. Cooperative theory explicitly identifies a ``user-owner'' principle: those who use a cooperative are also those who own and finance it \citep{dunn1988,zeuliradel2005}. Platform cooperativism later extended user and producer ownership into digital platforms; \citet{scholz2016} used the portmanteau ``produser'' for arrangements in which producer and user roles overlap, and \citet{schneider2018} analysed an ``Internet of ownership'' as an alternative to investor-only platform governance. These literatures make it inappropriate to claim that user ownership, participation-based ownership or online co-ownership originates here.

The present terminology has a more specific documented lineage. ERC-4950, created in March 2022 by Mu\~noz, de la Rosa and Easy Innova, explicitly uses the label ``Usowners'' for users who consume an NFT and become partial owners \citep{erc4950}. The proposal is a standards-track ERC whose current status is \emph{Stagnant}; it should therefore be treated as an open technical precursor, not an adopted Ethereum standard. \emph{Usownership} is the noun introduced here for the generalised institution: recurring, contribution-weighted, firm-specific equity acquisition inside firms that begin as predominantly conventional private corporations. It need not imply one-person-one-vote cooperative governance, full user financing, equal grants or immediate collective ownership. The novelty claim is consequently not that a user could become a partial owner, which ERC-4950 already stated, but that this user-to-owner relation is formalised as a scalable corporate ownership-transition mechanism and evaluated jointly with contribution measurement, dilution, income protection and monetary-financing constraints.

A bounded search conducted on 9 September 2026 also identified unrelated uses of the orthographic string \texttt{USOWNERSHIP}/\texttt{usOwnership} as a compact label for ``U.S. ownership'', including a contemporary financial-API field \citep{straitsx2026}. It did not identify an indexed source using \emph{usownership} as the noun for the generalised contribution-weighted institution defined here. The defensible lexical claim is therefore a new semantic noun in this research lineage, not ownership of an otherwise unused character string, proof of global lexical priority, or legal/trademark clearance. Search details are reported in the online supplement.

Wage-earner funds and collective capital institutions anticipated gradual ownership diffusion \citep{meidner1981,pontusson1992,furendal2024}. The AI Windfall Clause proposes advance sharing of exceptional AI gains \citep{windfall2020}. The closest contemporary fiscal antecedent is \citet{bearerfriend2025}, who propose an in-kind equity tax on generative-AI firms. Recurring grants to heterogeneous firm-specific participants differ from their one-time public-ownership instrument in incidence, beneficiaries and timing, not in the general discovery that AI equity can be redistributed.

Ownership is not merely the receipt of a labelled dividend. In incomplete-contract theory, residual control matters and its reallocation creates both benefits and distortions \citep{grossmanhart1986}. Our simulated shares give economic claims but do not endogenise voting, corporate strategy or managerial incentives. A governance advantage therefore remains a testable institutional hypothesis. Evidence from employee ownership is suggestive but not a transferable coefficient: the meta-analysis of \citet{oboyle2016} finds a small positive association with performance, not proof that user ownership raises productivity or that its dilution cost is zero.

User contributions have separate literatures. Data-as-labour proposals address compensation for data inputs \citep{dataaslabour2018}; \citet{jones2020data} study nonrival data, ownership and privacy incentives. Data Shapley provides task-specific valuation techniques \citep{ghorbani2019}; DataBright considers decentralised data ownership and rewards \citep{databright2018}. A prediction-accuracy contribution is not a universal measure of economic surplus, and a corporate share does not transfer data-processing consent. Usownership consequently needs a disclosed distributive rule, contractual boundaries and appeal procedures. It should not claim to measure each person's true marginal product merely by assigning a score.

\subsection{Monetary limits, tokenisation and the Internet of Value}
Seigniorage, inflation taxation and fiscal dominance are established subjects \citep{bailey1956,friedman1971,sargent1981,buiter2014}. Stock-flow discipline requires distinguishing issuance, cancellations, taxes, private dividends and public asset income \citep{godley2007}. Bank-created deposits are not sovereign money simply because they appear in a digital account \citep{mcleay2014}. The present sovereign-account closure excludes bank credit and legacy sovereign debt; it is deliberately not represented as a ready-made euro-area implementation.

Negative-rate proposals confront substitution into other stores of value \citep{agarwal2015}. CBDC studies analyse bank intermediation, payment efficiency and risks rather than presume that technology removes fiscal constraints \citep{auer2021,bidder2025,bindseil2025}. Circles already combines personal digital issuance and demurrage, although it is not a sovereign currency or a corporate-equity distribution system \citep{circles2026}. Neither demurrage nor digital basic income is a novelty claim here.

The implementation environment, however, is changing, and this paper also has a direct IoV precursor. De la Rosa Esteva's 2022 chapter on systemic risk in the Internet of Value argues that interoperable ledgers can support the democratisation of finance and property, analyses the fragility created by tokenised collateral and cross-ledger dependencies, and anticipates a ``massive migration of value onto the ledgers'' as a \emph{Value Deluge} requiring ownership and value-preservation mechanisms \citep{delarosa2022iov}. The chapter is conceptually relevant to the present migration-of-value premise, but it does not contain the present corporate contribution score, automation-linked issue rule or macroeconomic simulation. Those distinctions are important for a bounded novelty claim.

The term ``Internet of Value'' itself predates that chapter and this paper; Ripple was using it publicly by the middle of the 2010s and described it as an Internet on which money, securities and other assets can move as readily as information \citep{ripple2017}. The academic concept remains heterogeneous. \citet{lohmer2024} develop a definition around digital platforms, tokenisation and smart contracts, while the European Commission in 2026 explicitly described DLT and tokenisation as potentially paving the way to an Internet of Value \citep{eciov2026}. We use the expression descriptively, not as a prediction that all assets will migrate to one blockchain.

BIS work provides a more institutionally conservative formulation. Tokenisation records claims on real or financial assets on programmable platforms, integrating the asset record with rules governing transfer; programmable platforms can combine money and assets and support composability and atomic settlement \citep{bis2022,bis2023,bis2025,bis2026}. This is increasingly an operational rather than purely conceptual environment. By March 2026 the ECB reported close to EUR~4 billion of European DLT-based fixed-income issuance since 2021 and roughly EUR~1.6 billion of Eurosystem exploratory transactions across nine jurisdictions \citep{cipollone2026}. This is relevant to usownership because the proposed claim is repeated, granular and attached to corporate actions. Tokenisation can turn the right to receive, vest, transfer or distribute a participant share into an executable object. The economic claim still depends on company law, insolvency law, securities regulation and trusted information about the underlying firm.

``Programmable money'' also needs a boundary. The current digital-euro design explicitly rejects purpose-restricted money while allowing conditional payments \citep{ecbfaq2026}. Accordingly, the architecture studied here is better understood as \emph{programmable monetary infrastructure}: rules for issuance, cancellation, settlement and conditional transfer are machine-executable, but ordinary money is not restricted to a state-approved list of goods. This distinction limits both the surveillance risk and the claim of technological novelty.

\begin{table}[htbp]\centering\small
\caption{Closest mechanisms and the narrower question investigated}
\begin{tabularx}{\linewidth}{@{}p{.25\linewidth}p{.31\linewidth}X@{}}
\toprule
Literature or instrument & What is already established & Question added here \\
\midrule
Monetary finance & Income creation with money-demand and inflation constraints & How much of an ownership-transition gap fits inside that envelope?\\
Tax-transfer and robot taxes & Adaptive redistribution and transition insurance & Does equity improve anything beyond matched after-tax cash flows?\\
Industrial policy and public funds & Capacity support and publicly financed capital stocks & Can firm-specific participants acquire claims without an initial public purchase?\\
Wage-earner funds and AI equity tax & Gradual or one-off redistribution of productive claims & Recurring grants across a consumer--prosumer spectrum, initially unpooled\\
Data compensation and cooperatives & Participant value and alternative ownership institutions & Coverage, dilution, score integrity and monetary financing tested jointly\\
Internet of Value / tokenisation & Programmable assets, settlement and conditional execution & Whether granular usownership becomes cheaper to operate without changing its incidence\\
\bottomrule
\end{tabularx}
\end{table}
The contribution is an integrated comparison and a set of conditional mechanisms, not a priority claim for its individual ingredients. The review uses primary publications and official records available by 9 September 2026. Abstract-only access is recorded in the source ledger; no claims requiring unseen proofs are attributed to those sources. This is a focused, expanded review, not a database-exhaustive systematic review.
